\section{Υφιστάμενες Λύσεις και Αδυναμίες}

Στο οικοσύστημα των \ENG{SCORM Engines} υπάρχουν τόσο εμπορικές όσο και ανοιχτού κώδικα λύσεις. Παραδείγματα
εμπορικών πλατφορμών αποτελούν υπηρεσίες όπως το \ENG{Rustici SCORM Cloud}, το \ENG{Adobe Captivate Prime}
και το \ENG{iSpring Learn}. Από την πλευρά του ανοιχτού κώδικα, πλατφόρμες όπως το \ENG{Moodle} και το
\ENG{Ilias} παρέχουν ενσωματωμένη υποστήριξη για εκτέλεση \ENG{SCORM} περιεχομένου.

Παρά την ωριμότητα αυτών των συστημάτων, παρουσιάζονται ορισμένοι πρακτικοί περιορισμοί. Σε πολλές περιπτώσεις
ο μηχανισμός εκτέλεσης \ENG{SCORM} είναι στενά συνδεδεμένος με την εσωτερική αρχιτεκτονική του \ENG{LMS},
γεγονός που δυσχεραίνει την ανεξάρτητη επαναχρησιμοποίηση ή ενσωμάτωσή του σε άλλα πληροφοριακά συστήματα.

Επιπλέον, συχνά παρατηρείται περιορισμένη διαφάνεια στον τρόπο με τον οποίο διαχειρίζονται ζητήματα όπως η
ασφάλεια, ο έλεγχος πρόσβασης στα \ENG{runtime APIs} και η αποθήκευση της προόδου των
εκπαιδευομένων.

Άλλα προβλήματα σχετίζονται με τη διαχείριση κατάστασης και το λειτουργικό έλεγχο των συστημάτων.
Σε αρκετές υλοποιήσεις η πρόοδος αποθηκεύεται σε ένα ενιαίο επίπεδο χωρίς διαχωρισμό προσωρινής κατάστασης και
οριστικών δεδομένων, γεγονός που δυσκολεύει την ανθεκτικότητα και την ανάκτηση από σφάλματα. Παράλληλα,
η απουσία ολοκληρωμένων μηχανισμών \ENG{logging}, κεντρικής συλλογής αρχείων καταγραφής και βασικών
\ENG{health checks} καθιστά δυσκολότερη τη διάγνωση προβλημάτων σε πραγματικές αναπτύξεις.

Τέλος, ακόμη και όταν παρέχονται μορφές τεκμηρίωσης, συχνά απουσιάζουν συνοδευτικά εργαλεία όπως
\ENG{SDKs}, παραδείγματα διασύνδεσης και αναπαραγώγιμα περιβάλλοντα ανάπτυξης που θα διευκόλυναν την
ενσωμάτωση ενός τέτοιου εργαλείου σε τρίτες εφαρμογές.